% ---- Dataset ----------------------------------------------------------------
% Target length: ~1.5 pages.
% Structure: source, collection, license filter, normalization,
%            captioning pipeline, statistics, analysis.

\subsection{Source: The Iconify Ecosystem}

Iconify\footnote{\url{https://iconify.design}} is an open-source meta-library
aggregating over 300,000 icons from 225 icon sets. Unlike individual icon
libraries, Iconify provides a unified JSON format and a public GitHub repository
(\texttt{iconify/icon-sets}) containing all icons as structured data without
requiring web scraping. Each set carries explicit license metadata, making
per-set license filtering straightforward.

\subsection{License Filtering}
\label{sec:license}

Machine learning training use is governed by the license terms of the training
data. We apply a conservative license filter that retains only sets under
permissive open-source licenses and excludes all ShareAlike and copyleft
variants, as detailed in Table~\ref{tab:licenses}.

\begin{table}[h]
\centering
\caption{License filter applied to Iconify collections. ShareAlike and GPL
licenses are excluded due to their potential implications for model weight
distribution.}
\label{tab:licenses}
\begin{tabular}{lrrl}
\toprule
\textbf{License (SPDX)} & \textbf{Sets} & \textbf{Icons} & \textbf{Decision} \\
\midrule
MIT           & 99 & 114,830 & \textcolor{green!60!black}{Include} \\
Apache-2.0    & 30 &  82,644 & \textcolor{green!60!black}{Include} \\
CC-BY-4.0     & 51 &  55,648 & \textcolor{green!60!black}{Include} \\
CC0-1.0       &  8 &   8,051 & \textcolor{green!60!black}{Include} \\
OFL-1.1       & 13 &   6,211 & \textcolor{green!60!black}{Include} \\
CC-BY-3.0     &  1 &   4,123 & \textcolor{green!60!black}{Include} \\
ISC           &  3 &   2,998 & \textcolor{green!60!black}{Include} \\
CC-BY-NC-4.0  &  1 &     479 & \textcolor{green!60!black}{Include (non-commercial)} \\
Unlicense     &  1 &     430 & \textcolor{green!60!black}{Include} \\
MPL-2.0       &  1 &     288 & \textcolor{green!60!black}{Include} \\
BSD-3-Clause  &  1 &     210 & \textcolor{green!60!black}{Include} \\
\midrule
CC-BY-SA-4.0  &  7 &  21,828 & \textcolor{red}{Exclude (ShareAlike)} \\
CC-BY-NC-SA-4.0 & 1 & 1,577 & \textcolor{red}{Exclude (ShareAlike)} \\
CC-BY-SA-3.0  &  2 &      79 & \textcolor{red}{Exclude (ShareAlike)} \\
GPL-*         &  6 &   1,522 & \textcolor{red}{Exclude (Copyleft)} \\
\midrule
\textbf{Total included} & \textbf{209} & \textbf{275,912} & \\
\bottomrule
\end{tabular}
\end{table}

Additionally, 9,045 icons marked as \texttt{hidden} in the Iconify metadata
(deprecated or internal entries) are excluded, and alias entries (icons defined
as references to parent icons) are skipped to avoid duplicate training samples.

This study is conducted for non-commercial research purposes. Beyond license
terms, we note that (1) the EU DSM Directive (Article 4) explicitly permits
text and data mining of lawfully accessed content for research use; (2) the
trained model weights are a derivative work that does not reproduce training
icons verbatim. Each record in the released dataset retains the original
Iconify metadata fields \texttt{collection\_name}, \texttt{license\_spdx},
\texttt{author}, and \texttt{license\_url}, enabling downstream users to
comply with the attribution requirements of CC-BY and similar licenses.

\subsection{SVG Normalization}

Raw Iconify data consists of SVG \texttt{body} fragments — the inner content
of the \texttt{<svg>} element without the root tag or namespace declaration.
We assemble each icon into a complete, valid SVG document:

\begin{lstlisting}[style=svgstyle, caption={Assembled SVG document format}]
<svg xmlns="http://www.w3.org/2000/svg"
     viewBox="0 0 {W} {H}" width="{W}" height="{H}">
  {body}
</svg>
\end{lstlisting}

Two color conventions appear in the corpus. \emph{Monochrome} icons (the
majority) use \texttt{currentColor} as the fill and stroke value — the
standard convention in icon design that allows the icon to inherit its
context's text color. \emph{Multicolor} icons contain explicit color values
(hex codes, named colors, or \texttt{rgb()} expressions) hardcoded in the
path elements, e.g.\ \texttt{fill="\#4285F4"}. Each record carries an \texttt{is\_multicolor} boolean flag computed by
scanning the SVG body for explicit color values (hex codes, \texttt{rgb()},
or the named colors \texttt{white}/\texttt{black}) outside any
\texttt{currentColor} reference. This flag is used during curriculum
ordering (Section~\ref{sec:curriculum}) to expose the model to simple
monochrome icons before the more complex multicolor ones.

\subsection{Automated Caption Generation}
\label{sec:captioning}

Natural language descriptions are generated using Qwen2.5-VL-7B-Instruct
\citep{qwen25vl}, a vision language model running on Apple Silicon via the
MLX framework \citep{mlx2023}. Each icon is rendered to a 256$\times$256 PNG
using CairoSVG and passed to the VLM with a structured prompting template.

We generate two caption variants per icon:

\begin{itemize}
    \item \textbf{Short caption (5--12 words):} a natural search query expressing
          the concept and visual style, e.g., \textit{"outlined credit card with
          sparkles"} or \textit{"filled shopping cart with checkmark badge"}.
          This serves as the primary training prompt.
    \item \textbf{Long caption (20--35 words):} a detailed description including
          visual style (outline/filled/flat/solid/line), all visible elements,
          any decorations or modifiers, and likely UI context.
\end{itemize}

The icon name and collection name from Iconify metadata are provided as context
in the prompt, helping the model identify the semantic intent of ambiguous icons.
Caption generation runs checkpointed at the icon level, enabling resumable
processing across multiple sessions.

\todo{Report: total captioning time, model throughput, fraction of icons
requiring fallback to icon name (parse failure rate).}

\subsection{Dataset Statistics}

\todo{Add figures once captioning is complete:
(1) distribution of path\_count (histogram);
(2) proportion monochrome vs. multicolor by collection;
(3) caption length distribution (short and long);
(4) most frequent concepts in short captions (word cloud or bar chart);
(5) sample grid of icons across collections.}

Table~\ref{tab:dataset_stats} summarises the key characteristics of the
\dataset{} corpus.

\begin{table}[h]
\centering
\caption{\dataset{} corpus statistics.}
\label{tab:dataset_stats}
\begin{tabular}{lr}
\toprule
\textbf{Attribute} & \textbf{Value} \\
\midrule
Total icons            & 275,912 \\
Icon sets              & 209 \\
Monochrome icons       & \todo{X\%} \\
Multicolor icons       & \todo{X\%} \\
Median path count      & \todo{X} \\
Mean path count        & \todo{X} \\
Mean short caption length (words) & \todo{X} \\
Mean long caption length (words)  & \todo{X} \\
Training split         & 248,120 (90\%) \\
Validation split       &  13,896  (5\%) \\
Test split             &  13,896  (5\%) \\
\bottomrule
\end{tabular}
\end{table}
